% !TEX root = ../proyect.tex

% =====================================================
% APÉNDICES - placeholder
% =====================================================
% Material complementario que no encaja en el cuerpo principal pero conviene
% conservar como evidencia del trabajo realizado.
% Candidatos:
%   - Listado completo de casos de prueba (E2E, RAG, NLU).
%   - Esquema SQL de la base de datos.
%   - Configuración completa de Rasa (config.yml, domain.yml resumido).
%   - Informes de Clockify u otra herramienta de monitorización del tiempo
%     (recomendado por la guía §5).
%   - Capturas adicionales del panel admin y del widget.

\appendix
\chapter{Apéndices}\label{cap:apendices}

% =====================================================
\section{Actas de reuniones con el tutor}\label{ap:actas}
% =====================================================

Las reuniones de seguimiento con el tutor se celebraron al inicio de cada sprint. A continuación se recogen las actas en orden cronológico inverso.

\noindent\textbf{Datos de contacto}

\begin{tabular}{ll}
Tutor: & José Antonio Parejo Maestre, \texttt{japarejo@us.es} \\
Alumno: & Santia Bregu, \texttt{sanbre@alum.us.es} \\
\end{tabular}

% --------------------------------------------------
\subsection*{Sprint 8 (29/04/2026)}
% --------------------------------------------------

\textbf{Desarrollo de la reunión:}
\begin{itemize}
  \item Ajustar enfoque del resumen y añadir resumen en inglés.
  \item Incluir apéndice con actas de reuniones.
  \item Mejoras en manuales: capturas de pantalla, versiones de imágenes Docker y ejemplos de código.
  \item Revisión de diagramas y EDT.
  \item Pendiente: sección de metodología y revisión del capítulo de pruebas.
  \item Explicación de entregables y siguientes pasos del cierre del proyecto.
\end{itemize}

\textbf{Objetivos acordados:}
\begin{itemize}
  \item Reescribir el resumen y elaborar la sección de metodología.
  \item Añadir apéndice de actas y mejorar los manuales.
  \item Revisar el capítulo de pruebas.
  \item Enviar nueva versión de la memoria al tutor.
  \item Preparar el vídeo demostración del prototipo.
\end{itemize}

% --------------------------------------------------
\subsection*{Sprint 7 (15/04/2026)}
% --------------------------------------------------

\textbf{Desarrollo de la reunión:}
\begin{itemize}
  \item Enfoque en validación del sistema y escritura de la memoria.
  \item Definición de pruebas de aceptación (conversaciones e importaciones).
  \item Dos fases de validación: prototipo inicial y validación final.
  \item Uso de la guía de referencia para estructurar la memoria; secciones clave: introducción, conclusiones, manuales y validación.
  \item Tipos de manuales a redactar: usuario, despliegue y desarrollo.
  \item Envío progresivo de la memoria para obtener feedback continuo.
  \item Feedback positivo del estado actual del sistema.
  \item Planificación de sprint corto y fecha de próxima reunión.
\end{itemize}

\textbf{Objetivos acordados:}
\begin{itemize}
  \item Comenzar la redacción de la memoria con las secciones prioritarias.
  \item Definir y ejecutar el plan de pruebas de aceptación.
  \item Cerrar el desarrollo progresivamente y estabilizar el sistema.
\end{itemize}

% --------------------------------------------------
\subsection*{Sprint 6 (25/03/2026)}
% --------------------------------------------------

\textbf{Desarrollo de la reunión:}
\begin{itemize}
  \item Revisión de objetivos del sprint anterior (testing, memoria, sistema).
  \item Presentación del testing automatizado y manual; corrección de errores detectados.
  \item Mejoras en el sistema: mapeo de abreviaturas, \textit{fuzzy matching} y fallback con IA.
  \item Gestión de contexto con \textit{slots} para consultas encadenadas.
  \item Revisión del frontend: añadir enlaces a la web oficial de la US.
  \item Revisión de la memoria: EDT, arquitectura y decisiones de diseño.
  \item Correcciones en diagramas: uso de UML estándar.
  \item Revisión de horas dedicadas y mejora del tracking.
\end{itemize}

\textbf{Objetivos acordados:}
\begin{itemize}
  \item Desplegar el sistema y probarlo con usuarios reales.
  \item Dockerizar la aplicación.
  \item Implementar la épica de horarios y aulas.
  \item Mejorar la memoria: EDT y diagramas UML.
  \item Añadir enlaces en el frontend y aumentar horas dedicadas.
\end{itemize}

% --------------------------------------------------
\subsection*{Sprint 5 (27/02/2026)}
% --------------------------------------------------

\textbf{Desarrollo de la reunión:}
\begin{itemize}
  \item Revisión del estado actual del proyecto e identificación de tareas pendientes.
  \item Decisión de centrar el sprint en testing y consolidación del sistema.
  \item Definición del enfoque de pruebas: manual y automatizado.
  \item Propuesta de tareas opcionales: frontend y despliegue.
  \item Planificación de la próxima reunión.
\end{itemize}

\textbf{Objetivos acordados:}
\begin{itemize}
  \item Realizar testing exhaustivo: definir plan, ejecutar iteraciones y evaluar resultados.
  \item Identificar y corregir errores del sistema actual.
  \item Avanzar la memoria: arquitectura, decisiones de diseño y diagramas.
  \item Elaborar la EDT a nivel abstracto.
  \item Opcional: despliegue del sistema y versión inicial del frontend.
\end{itemize}

% --------------------------------------------------
\subsection*{Sprint 4 (03/02/2026)}
% --------------------------------------------------

\textbf{Desarrollo de la reunión:}
\begin{itemize}
  \item Revisión del avance de la épica de asignaturas.
  \item Planificación de la épica 7.3 (Asignaturas avanzado).
  \item Revisión de la estructura de la BD y del marco tecnológico.
\end{itemize}

\textbf{Objetivos acordados:}
\begin{itemize}
  \item Comenzar el desarrollo de la épica 7.3 (Asignaturas avanzado).
  \item Documentar en la plantilla TFG lo realizado: estructura de la BD, marco tecnológico y metodología.
  \item Mejorar las consultas actuales ampliando el contexto conversacional con el framework.
  \item Ampliar el modelo de información para dar soporte a múltiples titulaciones.
  \item Ampliar el conjunto de datos del módulo RAG y realizar pruebas.
\end{itemize}

% --------------------------------------------------
\subsection*{Sprint 3 (08/01/2026)}
% --------------------------------------------------

\textbf{Desarrollo de la reunión:}
\begin{itemize}
  \item Resumen del estado actual: datos de una titulación cargados (IS, ETSII) y frontend integrado con el backend Rasa.
\end{itemize}

\textbf{Objetivos acordados:}
\begin{itemize}
  \item Definir e implementar las historias de asignaturas: listado por titulación, créditos, cuatrimestre de pertenencia e impartición.
  \item Implementar pruebas de las historias anteriores.
  \item Comenzar la épica 7.3 (Asignaturas avanzado).
  \item Visualizar los vídeos del tutorial de integración de IA con Rasa.
  \item Documentar en la plantilla TFG: estructura de la BD, marco tecnológico y metodología.
  \item Comprobar infraestructura para ejecutar Ollama (modelo ligero tipo \texttt{llama3.2}).
\end{itemize}

% --------------------------------------------------
\subsection*{Sprint 2 (22/09/2025)}
% --------------------------------------------------

\textbf{Desarrollo de la reunión:}
\begin{itemize}
  \item Revisión del anteproyecto.
  \item Definición del alcance del proyecto.
  \item Fijación de objetivos para el sprint 2.
  \item Decisión: añadir un diccionario global de términos relevantes durante el desarrollo.
\end{itemize}

\textbf{Objetivos acordados:}
\begin{itemize}
  \item Finalizar el anteproyecto.
  \item Comenzar con el ítem del backlog ``Asignaturas''.
\end{itemize}

% --------------------------------------------------
\subsection*{Sprint 1 (28/07/2025)}
% --------------------------------------------------

\textbf{Desarrollo de la reunión:}
\begin{itemize}
  \item Revisión del anteproyecto.
  \item Creación de la infraestructura básica de trabajo.
  \item Adjudicación oficial del TFG con título provisional: \textit{``Diseño y desarrollo de un chat inteligente para consulta de información académica''}.
  \item Decisiones: usar GitHub como repositorio principal y registrar todas las horas desde el inicio, incluidas reuniones.
\end{itemize}

\textbf{Objetivos acordados:}
\begin{itemize}
  \item Justificar la elección tecnológica con estudio de alternativas y tabla comparativa.
  \item Montar la infraestructura: repositorio en GitHub y producto mínimo viable funcionando.
  \item Definir metodología y estándares de desarrollo.
  \item Elaborar la planificación inicial con horas aproximadas y objetivos por sprint.
\end{itemize}

\section{Informe de tiempo invertido}\label{ap:tiempo}

El registro horario del proyecto se ha llevado de forma continua mediante la herramienta \textbf{Clockify}, siguiendo la convención de etiquetado descrita en el Capítulo~\ref{cap:metodologia} (\texttt{SX} seguido de la descripción de la tarea, donde \texttt{X} es el número de sprint). El informe completo exportado desde Clockify, con una entrada por tarea registrada y su duración exacta, se adjunta al material del proyecto en formato CSV (\texttt{Clockify\_Time\_Report\_Summary\_01\_06\_2025-31\_05\_2026.csv}).

A modo de evidencia resumida, la Tabla~\ref{tab:clockify-categorias} muestra la distribución del tiempo por categoría de actividad y la Tabla~\ref{tab:clockify-sprints} reproduce el agregado por sprint, ambas calculadas a partir del fichero anterior con corte a fecha de \textbf{02/05/2026}.

\begin{table}[hbtp]
\centering
\small
\begin{tabular}{|l|r|r|}
\hline
\textbf{Categoría} & \textbf{Horas} & \textbf{\%} \\ \hline
Desarrollo            & 198,08 & 70,3\,\% \\ \hline
Documentación         &  26,89 &  9,5\,\% \\ \hline
Infraestructura       &  13,45 &  4,8\,\% \\ \hline
Investigación         &   8,73 &  3,1\,\% \\ \hline
Reuniones             &   5,67 &  2,0\,\% \\ \hline
Otras (anteproyecto, planning, fixes) & 28,90 & 10,3\,\% \\ \hline
\textbf{Total registrado} & \textbf{281,72} & \textbf{100,0\,\%} \\ \hline
\end{tabular}
\caption{Distribución del tiempo por categoría de actividad (Clockify, 02/05/2026).}
\label{tab:clockify-categorias}
\end{table}

\begin{table}[hbtp]
\centering
\small
\begin{tabular}{|c|r|}
\hline
\textbf{Sprint} & \textbf{Horas} \\ \hline
S1 (Anteproyecto)        &  20,16 \\ \hline
S2 (Infraestructura)     &  13,58 \\ \hline
S3 (Asignaturas)         &  33,19 \\ \hline
S4 (Asignaturas avanzado)&   8,17 \\ \hline
S5 (RAG y Gemini)        &  17,68 \\ \hline
S6 (Horarios y profesores)& 23,72 \\ \hline
S7 (Sprint largo: épicas, panel admin, piloto) & 150,52 \\ \hline
S8 (Cierre y memoria)    &  10,79 \\ \hline
Reuniones (sin sprint)   &   3,91 \\ \hline
\textbf{Total}           & \textbf{281,72} \\ \hline
\end{tabular}
\caption{Horas registradas por sprint (Clockify, 02/05/2026).}
\label{tab:clockify-sprints}
\end{table}

A esta cifra restan por contabilizar las horas correspondientes al cierre completo de la memoria (revisión final de manuales y capítulos pendientes), estimadas en aproximadamente 20\,h, que sitúan la previsión final del proyecto en torno a las 300\,h y por tanto dentro del margen establecido por la normativa.
